\documentclass[11pt]{article}
\usepackage[margin=1in]{geometry}
\usepackage{amsmath,amssymb,amsthm}
\usepackage[T1]{fontenc}
\usepackage{lmodern}
\usepackage{xcolor}
\usepackage{listings}
\usepackage{booktabs}
\usepackage{longtable}
\usepackage{array}
\usepackage[hidelinks]{hyperref}
\usepackage{url}

\newtheorem{theorem}{Theorem}
\newtheorem{lemma}[theorem]{Lemma}
\newtheorem{corollary}[theorem]{Corollary}
\newtheorem{proposition}[theorem]{Proposition}
\theoremstyle{definition}
\newtheorem{example}[theorem]{Example}
\theoremstyle{remark}
\newtheorem{remark}[theorem]{Remark}

% Code-like text (not checked by the name gate).
\newcommand{\lean}[1]{\texttt{#1}}
% A declaration of this development. Every \decl{...} is checked by the name gate
% (grep for a declaration of that name in CayleyBacharach/*.lean).
\DeclareUrlCommand\decl{\urlstyle{tt}\def\UrlBreaks{\do\_\do\.}}
\newcommand{\F}{\mathbb{F}}
\newcommand{\PP}{\mathbb{P}}
% version-2 commits: cayley-bacharach=d9c674c68c0a1b475a284f9bd33ec5f7e6d52689
\newcommand{\repo}{https://github.com/jpbald93/cayley-bacharach}
\newcommand{\commit}{9a4f11e82c291bc9428ad9b66bfaba388acf63de}
\newcommand{\shortcommit}{9a4f11e}
% \src{File}{n}: link to line n of CayleyBacharach/File.lean at \commit.
\newcommand{\src}[2]{\href{\repo/blob/\commit/CayleyBacharach/#1.lean\#L#2}{\texttt{#1:#2}}}

\lstdefinelanguage{lean}{
  morekeywords={theorem,lemma,def,class,example,by,fun,have,show,exact,intro,namespace,end,import,private,variable},
  sensitive=true, morecomment=[l]{--}, morecomment=[s]{/-}{-/}, morestring=[b]"}
\lstset{language=lean, basicstyle=\ttfamily\small, keywordstyle=\bfseries,
  columns=fullflexible, keepspaces=true, breaklines=true, frame=single, framerule=0.3pt,
  literate={≠}{{$\neq$}}1 {↔}{{$\leftrightarrow$}}2 {→}{{$\to$}}1
           {⨯₃}{{$\times_3$}}2 {⬝ᵥ}{{$\cdot_v$}}2 {₁}{{$_1$}}1 {₂}{{$_2$}}1 {₃}{{$_3$}}1
           {₄}{{$_4$}}1 {₅}{{$_5$}}1 {₆}{{$_6$}}1 {ℓ}{{$\ell$}}1 {≤}{{$\le$}}1
           {∀}{{$\forall$}}1 {∃}{{$\exists$}}1 {∧}{{$\wedge$}}1 {∨}{{$\vee$}}1
           {¬}{{$\neg$}}1 {∈}{{$\in$}}1 {ι}{{$\iota$}}1 {'}{{$^{\prime}$}}1}

\title{The Cayley--Bacharach theorem for plane cubics,\\ formalised in Lean~4,
with finite-field thresholds}
\author{Josh Bald\\ \small Independent researcher, Ontario, Canada\\ \small\texttt{jpbald93@gmail.com}}
\date{4 October 2026 (version 4)\\ \small Version~5, October 2026: adds Layers~14 and~15 (Section~\ref{sec:l1415})}

\emergencystretch=3em
\begin{document}
\maketitle

\begin{abstract}
We formalise the Cayley--Bacharach theorem for plane cubics in Lean~4 with Mathlib: under hypotheses made precise below, if two
cubics meet in nine points, every cubic through eight of them passes through the ninth. We prove
three versions. The first assumes that the eight points impose independent conditions on cubics;
we prove the exact criterion for this, in both directions: no five of the points are collinear
and the eight do not lie on a conic. The second assumes that the two cubics have no common line
or conic component. The third is the textbook form, in which the two cubics have exactly nine
common zeros. None of the three assumes an algebraically closed field. The first two hold over
every field. The third holds over every infinite field, and a finite field $K$ exactly when
$|K|\notin\{4,5,7\}$: we prove it for every other finite field, and it is false over $\F_4$,
$\F_5$ and $\F_7$, for which we give explicit counterexamples, checked by computer but not
formalised.
The line lemma behind all three (a cubic vanishing on a line contains it) needs a field with at
least three elements, and we formalise a counterexample over $\F_2$. As applications we derive
Pascal's and Pappus's theorems from Cayley--Bacharach under their natural incidence and
distinctness hypotheses, for fields with more than ten elements; version~5 adds direct proofs,
not through Cayley--Bacharach, of Pascal's and Pappus's theorems and of the Braikenridge--Maclaurin
converse over every field (Section~\ref{sec:l1415}). We also
prove that a conic is a line pair if and only if its Gram determinant vanishes, over an
algebraically closed field of characteristic different from two. The development has about
4{,}150 lines (the Layer-14/15 check files added in version~5 have about 2{,}100 more) and uses no \lean{sorry} and no axioms beyond Lean's standard three. The
mathematics is classical. We found no earlier formalisation of the theorem in the sources we
searched (Section~\ref{sec:related}).
\end{abstract}

\section{Introduction}

Let $K$ be a field. A \emph{cubic form} is a homogeneous polynomial of degree three in $x,y,z$ over $K$ (zero
included, so that the cubic forms make a vector space); a nonzero one defines a plane cubic. The following statement goes back to Cayley~\cite{Cayley} and
Bacharach~\cite{Bacharach}. Eisenbud, Green and Harris~\cite{EGH} survey its history and its
many generalisations.

\begin{theorem}[Cayley--Bacharach, textbook form]\label{thm:textbook}
Let $K$ be algebraically closed and let $F,G$ be plane cubics that meet in exactly nine distinct points $P_1,\dots,P_9$. If a cubic
$H$ passes through $P_1,\dots,P_8$, then it passes through $P_9$.
\end{theorem}

This is the classical statement, over an algebraically closed field; the versions we formalise are
stated precisely in Sections~\ref{sec:criterion}--\ref{sec:apps}. Textbooks also count intersection points with
multiplicity. Fulton~\cite[\S5.6]{Fulton}, for instance, works over an algebraically closed field
and derives Pascal's and Pappus's theorems from Max Noether's fundamental theorem. We make no
such assumption. Instead,
we work with the $K$-rational points of $\PP^2(K)$ and ask exactly which hypotheses make the
theorem true.

The proof we formalise is the linear-algebra proof. The cubics form a vector space of
dimension~10. Each point imposes one linear condition, so the cubics through eight points form a
space of dimension at least~2. If the dimension is exactly~2, then this space is spanned by $F$
and $G$, and every cubic in it vanishes at~$P_9$. Everything therefore depends on knowing when
eight points impose independent conditions on cubics. We prove the exact criterion
(Theorem~\ref{thm:criterion}): this happens if and only if no five of the points are collinear
and the eight points do not lie on a conic. A commonly quoted condition is ``no four collinear
and no seven on a nondegenerate conic'' (see e.g.~\cite{Wikipedia}). That condition is
sufficient, but it is not necessary, as the formal theorem and the explicit configurations of
Example~\ref{ex:criterion} show.

\paragraph{Contributions.}
\begin{itemize}
\item Both directions of the eight-point criterion over every field
(Theorem~\ref{thm:criterion}).
\item Cayley--Bacharach in three forms: general position (Theorem~\ref{thm:gp}), no common
component (Theorem~\ref{thm:ncf}), and exactly nine common zeros (Theorem~\ref{thm:classical}).
The first two hold over every field. The third holds over every infinite field and every finite
field with $|K|\notin\{4,5,7\}$.
\item Sharp field-size results (Section~\ref{sec:finite}). The line lemma holds if and only if
$|K|\ge3$, and the $\F_2$ counterexample is formalised. The textbook form holds over $\F_2$ (vacuously), $\F_3$ and every
finite field with at least eight elements, all formalised, and fails over $\F_4$, $\F_5$ and
$\F_7$. We give explicit counterexamples for these three, checked by computer but not formalised.
\item Pascal's and Pappus's theorems derived from Cayley--Bacharach with only their natural
incidence and distinctness hypotheses (for $|K|>10$), and the Gram-determinant criterion for line-pair conics (Section~\ref{sec:apps}).
Version~5 adds direct proofs of Pascal, Pappus and Braikenridge--Maclaurin over every field
(Section~\ref{sec:l1415}).
\item Version~5: the main theorems restated with Mathlib-only types, each with an explicit
non-degenerate witness, and a comparison with Lean statements written blind from textbook statements
(Section~\ref{sec:l1415}).
\item Lemmas that help a reader check that the formal statements say what they should
(Section~\ref{sec:correct}).
\end{itemize}

\paragraph{Availability.}
The code is at \url{\repo}. Line references in this paper are to commit \texttt{\shortcommit}.
The toolchain is Lean \texttt{v4.33.1}~\cite{Lean4} with Mathlib~\cite{Mathlib} pinned in
\texttt{lake-manifest.json}. Apart from the type abbreviation \decl{PlaneCubic}, all declarations of the main files are in the namespace \lean{PlaneCubic}; the version-5 files use the namespaces \lean{CBL14} and \lean{CayleyBacharach.Classical} (Section~\ref{sec:l1415}).
Table~\ref{tab:thms} lists the main formal statements.
Version~5 of this paper adds the Layer-14 and Layer-15 files of Section~\ref{sec:l1415} in a new
commit, \texttt{d9c674c68c0a1b475a284f9bd33ec5f7e6d52689}. The earlier statements and library files are unchanged in that commit
(only the root file gains two import lines, and \path{gate.sh}, \path{.gitignore} and the README are
updated), so the line references above are valid at both commits.
Supporting material (the Layer-14 and Layer-15 reports, the blind writer's statements and report, the transcript audit, and the reviews of this version with the fixes made in response) is in the folder \texttt{Cayley-Bacharach} of the repository \url{https://github.com/jpbald93/LEAN-verification}.

{\small
\setlength{\LTcapwidth}{\textwidth}
\begin{longtable}{@{}>{\raggedright\arraybackslash}p{0.40\textwidth}
  >{\raggedright\arraybackslash}p{0.33\textwidth}l@{}}
\caption{Informal statements and their formal counterparts. The field $K$ is arbitrary unless a
hypothesis is shown. Line numbers refer to the first line of the declaration.}\label{tab:thms}\\
\toprule
Informal statement & Lean name & Source \\ \midrule
\endfirsthead
\toprule
Informal statement & Lean name & Source \\ \midrule
\endhead
\bottomrule
\endfoot
\multicolumn{3}{@{}l}{\emph{Representation and meaning}}\\
Cubics as coefficient vectors & \decl{PlaneCubic} & \src{Defs}{16} \\
Vanishing does not depend on the representative & \decl{vanishesAt_smul_iff} & \src{Defs}{86} \\
Evaluation of a product $\ell\cdot q$ & \decl{eval_mulLin} & \src{Line}{32} \\
Evaluation of a line pair $\ell\cdot m$ & \decl{evalC_linMul} & \src{Conic}{40} \\
Kernel of evaluation $=$ cubics through the points & \decl{ker_evalMap_eq} & \src{Through}{35} \\
Common zero of $F$ and $G$ & \decl{CommonZero} & \src{Classical}{231} \\
The common zeros are exactly the $P_i$ & \decl{ExactlyMeetIn} & \src{Classical}{235} \\
\multicolumn{3}{@{}l}{\emph{Lines and conics}}\\
Line lemma, $|K|\ge 3$ (Lemma~\ref{lem:line}) & \decl{line_lemma} & \src{Line}{247} \\
Line lemma, finite $K$, $3\le|K|$ & \decl{line_lemma_fintype} & \src{FiniteField}{34} \\
Line lemma fails over $\F_2$ & \decl{line_lemma_false_zmod_two} & \src{Line}{270} \\
Conic line lemma, every field & \decl{conic_line_lemma} & \src{Conic}{167} \\
Three collinear points on a conic & \decl{conic_three} & \src{LineRestrict}{105} \\
Four collinear points on a cubic, $|K|\ge3$ & \decl{cubic_four} & \src{LineRestrict}{124} \\
\multicolumn{3}{@{}l}{\emph{The eight-point criterion (Theorem~\ref{thm:criterion})}}\\
General position $\Rightarrow$ dimension $\le 2$, $|K|\ge3$ & \decl{finrank_cubicsThrough_le_two} & \src{General}{185} \\
General position $\Rightarrow$ dimension $=2$, $|K|\ge3$ & \decl{finrank_ker_evalMap_eq_two} & \src{General}{209} \\
General position $\Rightarrow$ dimension $=2$, every field & \decl{finrank_ker_evalMap_eq_two_any} & \src{FiniteField}{92} \\
Five collinear $\Rightarrow$ dimension $\ge 3$ & \decl{finrank_ker_ge_of_line} & \src{Criterion}{122} \\
On a conic $\Rightarrow$ dimension $\ge 3$ & \decl{finrank_ker_ge_of_conic} & \src{Criterion}{102} \\
The criterion as a single equivalence, every field & \decl{finrank_ker_evalMap_eq_two_iff} & \src{CriterionIff}{30} \\
Over $\F_2$ at most seven distinct points & \decl{card_le_seven_of_not_atLeastThree} & \src{FiniteField}{55} \\
\multicolumn{3}{@{}l}{\emph{Cayley--Bacharach}}\\
Cubics through the eight $=\operatorname{span}\{F,G\}$, $|K|\ge3$ & \decl{mem_span_of_eight} & \src{Theorem}{19} \\
General position, $|K|\ge3$ & \decl{cayley_bacharach} & \src{Theorem}{48} \\
General position, every field (Thm.~\ref{thm:gp}) & \decl{cayley_bacharach_any} & \src{FiniteField}{103} \\
No common factor $\Rightarrow$ general position, $|K|\ge3$ & \decl{general_position_of_no_common_factor} & \src{Factor}{98} \\
No common factor, $|K|\ge3$ & \decl{cayley_bacharach_of_no_common_factor} & \src{Factor}{158} \\
No common factor, every field (Thm.~\ref{thm:ncf}) & \decl{cayley_bacharach_of_no_common_factor_any} & \src{FiniteField}{117} \\
Exactly nine, infinite $K$ (Thm.~\ref{thm:classical}) & \decl{cayley_bacharach_classical} & \src{Classical}{242} \\
Exactly nine, finite $K$, $|K|\ge 11$ & \decl{cayley_bacharach_classical_fintype} & \src{ClassicalFinite}{203} \\
Exactly nine, finite $K$, $|K|\ge 8$ (Thm.~\ref{thm:classical}) & \decl{cayley_bacharach_classical_card8} & \src{ClassicalSmall}{316} \\
Nine common zeros, $|K|=3$ (Thm.~\ref{thm:f3}) & \decl{cayley_bacharach_card3} & \src{ClassicalF3}{231} \\
Exactly nine, $|K|=3$ & \decl{cayley_bacharach_classical_card3} & \src{ClassicalF3}{243} \\
Exactly nine, $K=\F_3$ & \decl{cayley_bacharach_classical_zmod3} & \src{ClassicalF3}{255} \\
Exactly nine, finite $K$, $|K|\notin\{4,5,7\}$ (Thm.~\ref{thm:f3}) & \decl{cayley_bacharach_classical_finite} & \src{ClassicalF3}{266} \\
$|K|=3$: no common line & \decl{no_common_line_card3} & \src{ClassicalF3}{125} \\
$|K|=3$: no common conic & \decl{no_common_conic_card3} & \src{ClassicalF3}{188} \\
Construction of $|K|+1$ distinct points on a line, finite $K$ & \decl{line_points} & \src{ClassicalSmall}{42} \\
A conic with no line component and two zeros has $\ge|K|+1$ zeros & \decl{conic_points} & \src{ClassicalSmall}{157} \\
\multicolumn{3}{@{}l}{\emph{Applications (Section~\ref{sec:apps}); Pascal and Pappus need $|K|>10$}}\\
Pascal from Cayley--Bacharach, explicit genericity & \decl{pascal_of_cayley_bacharach} & \src{Pascal}{46} \\
Pappus from Cayley--Bacharach, explicit genericity & \decl{pappus_of_cayley_bacharach} & \src{Pascal}{179} \\
Pascal, natural hypotheses (Thm.~\ref{thm:pascal}) & \decl{pascal_clean} & \src{PascalClean}{331} \\
Pappus, natural hypotheses (Thm.~\ref{thm:pappus}) & \decl{pappus_clean} & \src{PappusClean}{178} \\
Line pair $\Rightarrow$ Gram determinant $0$, $2\ne0$ & \decl{det_gram_linMul} & \src{PappusClean}{223} \\
$2\ne0$ is needed & \decl{det_gram_linMul_char_two} & \src{PappusClean}{252} \\
Pascal with $\det M(q)\ne0$, $2\ne0$ & \decl{pascal_clean_det} & \src{PappusClean}{236} \\
Gram determinant $0$ $\Leftrightarrow$ line pair, $K$ alg.\ closed, $2\ne0$ & \decl{det_gram_eq_zero_iff} & \src{PappusClean}{316} \\
\end{longtable}
}

\section{Representation and conventions}\label{sec:conv}

A cubic is its vector of ten coefficients, \decl{PlaneCubic}~\lean{K := Fin 10} $\to$ \lean{K}, in the
monomial order $x^3,x^2y,x^2z,xy^2,xyz,xz^2,y^3,y^2z,yz^2,z^3$. The name avoids a clash with
Mathlib's \lean{Cubic}. Because it is an abbreviation of a function type, the vector-space
structure and \lean{Module.finrank K (PlaneCubic K) = 10} come directly from Mathlib. A conic is a
vector \lean{Fin 6} $\to$ \lean{K} in the order $x^2,xy,xz,y^2,yz,z^2$, and a linear form or line is a
vector $\ell$ \lean{: Fin 3} $\to$ \lean{K}. The functions \decl{eval} and \decl{evalC} evaluate cubics and
conics at a vector $v\in K^3$. The function \decl{mulLin}~$\ell$~$q$ is the cubic $\ell\cdot q$
and \decl{linMul}~$\ell$~$m$ is the conic $\ell\cdot m$; both are given by explicit coefficient
formulas.

A point of $\PP^2(K)$ is represented by a nonzero vector. We do not use Mathlib's
\lean{Projectivization} here. Instead, a family of points \lean{P : Fin n} $\to$ \lean{Fin 3} $\to$ \lean{K} is
required to be pairwise non-proportional:
\begin{center}
\lean{hP :} $\forall\, i\,j,\ i\ne j\to P_i\times P_j\neq 0$.
\end{center}
This says that the $P_i$ are distinct points of $\PP^2(K)$. It also forces each $P_i$ to be
nonzero when $n\ge 2$ (\decl{ne_zero_of_pairwise}). A point $v$ lies on the line $\ell$ when
$\ell\cdot v=0$, and the line through $u\ne v$ is $u\times v$. The two general-position
hypotheses on eight points read as follows:
\begin{align*}
\text{\lean{h5}}&:\ \forall\,\ell\ne0,\ \forall\, s\subseteq\{0,\dots,7\},\
  (\forall i\in s,\ \ell\cdot P_i=0)\to |s|\le4 &&\text{(no five collinear)},\\
\text{\lean{hconic}}&:\ \forall\, q,\ (\forall i,\ q(P_i)=0)\to q=0 &&\text{(not all on a conic)}.
\end{align*}
Here ``conic'' includes line pairs and double lines. The hypothesis \lean{hconic} is therefore the
same as saying that the eight evaluation functionals on conics span a space of dimension six.

Two typeclasses express field size. \decl{AtLeastThree}~\lean{K} says that $K$ has an element
other than $0$ and~$1$. There is an instance for every infinite field
(\decl{AtLeastThree.of_infinite}), and for a finite field it is equivalent to $3\le|K|$
(\decl{atLeastThree_iff_card}). \decl{CardGT}~\lean{K n} says that $K$ has a finite subset with
more than $n$ elements. It supplies the older counting interface ($n=10$) used by the Pascal and Pappus files and by a variant of the classical theorem in \texttt{ClassicalFinite.lean}.

\section{Lines, conics and the eight-point criterion}\label{sec:criterion}

\subsection{Restricting to a line}

\begin{lemma}[line lemma]\label{lem:line}
Let $|K|\ge3$ and let $\ell\ne0$. If a cubic $F$ vanishes at every point of the line $\ell=0$,
then $F=\ell\cdot q$ for a conic $q$. Over every field, a conic vanishing on the line
$\ell=0$ is $\ell\cdot m$ for a linear form~$m$.
\end{lemma}

\begin{proof}
Choose a coordinate $k$ with $\ell_k\ne0$ and parametrise the line with $x_k$ eliminated. The
restriction of $F$ is a binary cubic $a s^3+b s^2t+c st^2+dt^3$. At $[1:0]$ we obtain $a=0$. On the chart $[s:t]=[r:1]$, evaluating $a r^3+b r^2+c r+d$ at
$r=0,1,r_0$, with $r_0\notin\{0,1\}$, gives $b=c=d=0$. This is
\decl{coeffs_eq_zero_of_forall_three}. The quotient $q$ is then given by an explicit
formula, obtained by pseudo-division of $\ell_k^3F$ by $\ell$ in the variable~$x_k$, and each of
the ten coefficient identities is closed by \lean{linear\_combination}. For conics, every line
has the three points $[1:0]$, $r=0$ and $r=1$. These suffice for a binary quadratic, so no size
hypothesis is needed.
\end{proof}

The first statement is false over $\F_2$. The cubic $x^2y+xy^2$ vanishes at the three points
$(1:0:0)$, $(0:1:0)$, $(1:1:0)$ of $z=0$, but it is not divisible by~$z$. This is
\decl{line_lemma_false_zmod_two}. Two consequences are used throughout: a conic that vanishes at
three distinct points of a line contains the line (\decl{conic_three}, every field), and a cubic
that vanishes at four distinct points of a line contains the line (\decl{cubic_four}, $|K|\ge3$).

\subsection{When eight points impose independent conditions}

For points $P_1,\dots,P_8$, let $V(P)$ be the space of cubics vanishing at all of them. It is the
kernel of the evaluation map $K^{10}\to K^8$ (\decl{ker_evalMap_eq}), so $\dim V(P)\ge 2$.

\begin{theorem}[eight-point criterion]\label{thm:criterion}
Let $K$ be any field and let $P_1,\dots,P_8$ be distinct points of $\PP^2(K)$. Then
$\dim V(P)=2$ if and only if no five of the points are collinear and the eight points do not all
lie on a conic.
\end{theorem}

The formal version is stated both as three separate implications and as a single equivalence.
\decl{finrank_ker_evalMap_eq_two_any} gives
$\dim V(P)=2$ under \lean{hP}, \lean{h5} and \lean{hconic}. For the converse,
\decl{finrank_ker_ge_of_line} gives $\dim V(P)\ge3$ if all but at most three of the points lie on
a line, and \decl{finrank_ker_ge_of_conic} gives $\dim V(P)\ge 3$ if a nonzero conic passes
through all of them. The ``if and only if'' itself is
\decl{finrank_ker_evalMap_eq_two_iff}, which combines these three: $\dim V(P)=2$ if and only if
no line carries five of the points and no nonzero conic passes through all eight.

\begin{proof}[Proof of the hard direction]
Suppose $\dim V(P)\ge 3$. We use four facts about conics. Four distinct points that are not all
collinear lie on at most a $2$-dimensional space of conics (\decl{conics_le_two}). Five
distinct points, no four collinear, lie on at most a $1$-dimensional space of conics
(\decl{conics_le_one}). If $u\ne v$ are among the points, some nonzero conic vanishes at every
point off the line $uv$ (\decl{exists_conic_off_line}). If moreover a third point $w$ lies on
$uv$, the conics through the points off $uv$ form a space of dimension at least~$2$
(\decl{conics_off_line_ge_two}). The last two facts come from adding one or two further points on
$uv$ and applying \decl{cubic_four}. Choosing these points needs an element $t\notin\{0,1\}$,
which is where $|K|\ge3$ enters. There are three cases.

\emph{(A) Some line $\ell$ contains exactly four of the points.} Every cubic through the eight
contains~$\ell$, so $\dim V(P)$ is at most the dimension of the space of conics through the other
four (\decl{cubics_resid}). Those four are not collinear, since otherwise $\ell$ times their line
would be a conic through all eight. Hence $\dim V(P)\le 2$.

\emph{(B) No line contains four of the points, but three are collinear, on a line~$\ell$.} The
five points off $\ell$ lie on a space of conics of dimension at least~$2$, but no four of them
are collinear, so this dimension is at most~$1$.

\emph{(C) No three of the points are collinear.} For each pair $\{0,1\}$, $\{0,2\}$, $\{1,2\}$
take a nonzero conic through the six points off the line joining the pair. All three pass through
$P_3,\dots,P_7$, whose conics form a space of dimension at most~$1$. So the three conics are
proportional, and the first one vanishes at all eight points, a contradiction with \lean{hconic}.

This proves the theorem for $|K|\ge3$ (\decl{finrank_cubicsThrough_le_two},
\decl{finrank_ker_evalMap_eq_two}). Over $\F_2$ the plane has only seven points, so eight
distinct points do not exist (\decl{card_le_seven_of_not_atLeastThree}), and the statement is
vacuous.
\end{proof}

\begin{example}\label{ex:criterion}
The hypotheses ``no four collinear'' and ``no seven on a conic'' are not necessary. We computed
$\dim V(P)$ in exact rational arithmetic for the following configurations.
% Script: code/criterion_examples.py in the repository.
\begin{itemize}
\item $(1:0:1),(2:0:1),(3:0:1),(5:0:1)$ on the line $y=0$, together with $(1:2:3)$, $(4:1:7)$,
$(2:5:1)$, $(3:3:2)$: $\dim V(P)=2$.
\item $(1:t:t^2)$ for $t=0,\dots,6$, on the conic $y^2=xz$, together with $(1:1:5)$:
$\dim V(P)=2$.
\item $(1:0:1),(2:0:1),(3:0:1),(5:0:1),(7:0:1)$ on $y=0$, together with $(1:2:3)$, $(4:1:7)$,
$(2:5:1)$: $\dim V(P)=3$.
\item $(1:t:t^2)$ for $t=0,\dots,7$: $\dim V(P)=3$.
\end{itemize}
The first two configurations satisfy the hypotheses of \decl{finrank_ker_evalMap_eq_two_any},
which allows four points on a line and seven on a conic. The conditions ``no four collinear, no
seven on a nondegenerate conic'' imply ours: eight points on a line pair put at least four on
one of the lines. They are therefore sufficient, but they exclude valid configurations.
\end{example}

\section{Cayley--Bacharach}\label{sec:cb}

\subsection{General position}

\begin{theorem}\label{thm:gp}
Let $K$ be any field and $P_1,\dots,P_9$ distinct points of $\PP^2(K)$. Suppose no five of
$P_1,\dots,P_8$ are collinear and $P_1,\dots,P_8$ do not lie on a conic. Let $F,G$ be linearly
independent cubics vanishing at all nine points. Then every cubic $H$ vanishing at
$P_1,\dots,P_8$ vanishes at~$P_9$.
\end{theorem}

\begin{lstlisting}
theorem cayley_bacharach_any (P : Fin 9 → Fin 3 → K)
    (hP : ∀ i j, i ≠ j → P i ⨯₃ P j ≠ 0)
    (h5 : ∀ ℓ : Fin 3 → K, ℓ ≠ 0 → ∀ s : Finset (Fin 8),
      (∀ i ∈ s, ℓ ⬝ᵥ P (Fin.castSucc i) = 0) → s.card ≤ 4)
    (hconic : ∀ q : Fin 6 → K, (∀ i : Fin 8, evalC q (P (Fin.castSucc i)) = 0) → q = 0)
    (F G : PlaneCubic K) (hFG : LinearIndependent K ![F, G])
    (hF : ∀ i, eval F (P i) = 0) (hG : ∀ i, eval G (P i) = 0)
    (H : PlaneCubic K) (hH : ∀ i : Fin 8, eval H (P (Fin.castSucc i)) = 0) :
    eval H (P 8) = 0
\end{lstlisting}

\begin{proof}
$\operatorname{span}\{F,G\}$ has dimension~2 and lies in $V(P_1,\dots,P_8)$, which has dimension
at most~2 by Theorem~\ref{thm:criterion}. So the two spaces are equal
(\decl{mem_span_of_eight}), and $H$ is a combination of $F$ and~$G$.
\end{proof}

The indexing is zero-based: \lean{P 8} is the ninth point and \lean{Fin.castSucc} embeds the
first eight. Linear independence of $F$ and $G$ is needed. If $G=cF$, the cubics through the eight
points need not all pass through the ninth.

\subsection{No common component}

\begin{theorem}\label{thm:ncf}
Let $K$ be any field and $P_1,\dots,P_9$ distinct points. Let $F,G$ be linearly independent
cubics through all nine points with no common linear factor and no common conic factor. Then
every cubic through $P_1,\dots,P_8$ passes through $P_9$.
\end{theorem}

The formal hypotheses are
\begin{lstlisting}
(hlin : ¬ ∃ ℓ q q', ℓ ≠ 0 ∧ F = mulLin ℓ q ∧ G = mulLin ℓ q')
(hcon : ¬ ∃ q m m', q ≠ 0 ∧ F = mulLin m q ∧ G = mulLin m' q)
\end{lstlisting}
They are statements about factorisations over~$K$.

\begin{proof}
We show that $P_1,\dots,P_8$ are in general position (\decl{general_position_of_no_common_factor}).
We in fact get the stronger conclusion that no four of them are collinear. If four were on a line
$\ell$, both $F$ and $G$ would contain~$\ell$ by \decl{cubic_four}. Suppose the eight points lie
on a nonzero conic~$q$. If three of them are collinear, $q$ contains their line
(\decl{conic_three}), so $q$ is a line pair and one of the two lines carries at least four of the
points, which was just excluded. If no three are collinear, the cubics through the eight form a
space of dimension at most three (\decl{finrank_le_three_of_no3}). This space contains the
three-dimensional space $\{m\cdot q\}$, so the two are equal, and then $q$ divides both $F$ and~$G$.
Now apply Theorem~\ref{thm:gp}.
\end{proof}

\subsection{Exactly nine common zeros}

The textbook hypothesis is that the common zeros of $F$ and $G$ are exactly the nine points. We
formalise it literally, for $K$-rational points:
\begin{lstlisting}
def CommonZero (F G : PlaneCubic K) (v : Fin 3 → K) : Prop :=
  v ≠ 0 ∧ eval F v = 0 ∧ eval G v = 0
def ExactlyMeetIn {ι : Type*} (F G : PlaneCubic K) (P : ι → Fin 3 → K) : Prop :=
  (∀ i, CommonZero F G (P i)) ∧ ∀ v, CommonZero F G v → ∃ i, v ⨯₃ P i = 0
\end{lstlisting}

\begin{theorem}\label{thm:classical}
Let $K$ be an infinite field, or a finite field with $|K|\ge8$. Let $P_1,\dots,P_9$ be distinct
points of $\PP^2(K)$ and let $F,G$ be linearly independent cubics whose common zeros in
$\PP^2(K)$ are exactly $P_1,\dots,P_9$. Then every cubic through $P_1,\dots,P_8$ passes
through~$P_9$.
\end{theorem}

The infinite case is \decl{cayley_bacharach_classical}; the finite case is
\decl{cayley_bacharach_classical_card8}. Both have the same hypotheses as
Theorem~\ref{thm:ncf}, with \lean{hex : ExactlyMeetIn F G P} in place of the factor hypotheses.

\begin{proof}[Proof for infinite $K$]
We show that exactness rules out a common factor, and then apply Theorem~\ref{thm:ncf}. A common
line would contribute infinitely many common zeros (\decl{no_common_line_of_exact}). Suppose
$F=m\,q$ and $G=m'q$. If $q$ vanishes at two of the $P_i$, it has infinitely many pairwise
non-proportional zeros (\decl{false_of_conic}), and all of them are common zeros. Otherwise $q$
vanishes at most at one of $P_1,P_2,P_3$. Then $m$ and $m'$ both vanish at two of these points,
so both are multiples of the line through them, and that line consists of common zeros
(\decl{no_common_conic_of_exact}).
\end{proof}

The step through \decl{false_of_conic} is where algebraic closure might seem necessary. Over
$\mathbb{R}$ the conic $x^2+y^2+z^2$ has no zeros, so a common conic factor does not by itself
produce infinitely many common zeros. The proof avoids the issue because the nine points are
themselves common zeros: a conic with two non-proportional zeros $u,v$ over an infinite field has
infinitely many. Projecting from~$u$, each $d$ gives the zero $q(d)\,u-B(u,d)\,d$, where $B$ is
the polar form. This works in every characteristic. If $B(u,v)=0$, the whole line $uv$ lies on
the conic instead.

\begin{proof}[Proof for finite $K$ with $|K|\ge 8$]
The same argument needs point counts. A line has exactly $|K|+1$ points; the proof needs only the lower bound, and \decl{line_points} constructs $|K|+1$ pairwise non-proportional points on it. A
conic with two distinct zeros and no line component has at least $|K|+1$ zeros
(\decl{conic_points}). For this, project from a zero $u$ onto all $|K|+1$ lines through~$u$.
Each line through $u$ contains a zero other than~$u$, except the tangent line, where the
construction returns $u$ itself. At most one line through $u$ is tangent, since two tangent
directions would force a line component (\decl{line_comp_of_polar}). Now suppose $F$ and $G$ have a
common line, or a common conic without line component through two of the~$P_i$. Every point of it
is a common zero, hence one of the~$P_i$. It has at least $|K|+1\ge9$ distinct points, so all
nine $P_i$ lie on it (\decl{card_add_one_le}). If they lie on a line, a cubic through eight of them
contains the line (\decl{classical_line_case}). If they lie on such a conic, no three of the
points are collinear, and a cubic through eight of them is a multiple of the conic
(\decl{classical_conic_case}). In either case $H(P_9)=0$. A common line-pair factor contains a
common line. If a common conic factor vanishes at at most one of $P_1,P_2,P_3$, then, as in the
infinite case, $F$ and $G$ share the line through two of these points, and the line case applies.
If there is no common factor, Theorem~\ref{thm:ncf} applies.
\end{proof}

An earlier, cruder count, $|K|-1$ zeros on the conic instead of $|K|+1$, gave the result only for
$|K|\ge 11$ (\decl{cayley_bacharach_classical_card}, \decl{cayley_bacharach_classical_fintype}).
These versions are kept; the Pascal and Pappus files use them.

\begin{remark}[meaning over non-closed fields]
\lean{ExactlyMeetIn} refers to $K$-rational points only. Over $\mathbb{R}$, a pencil some of
whose nine complex base points are not real does not satisfy \lean{ExactlyMeetIn} for any nine
real points, so Theorem~\ref{thm:classical} says nothing about it. Over an algebraically closed
field, ``exactly nine distinct common zeros'' is the usual hypothesis. It excludes a common
component, so every intersection multiplicity is~1 by B\'ezout, and
Theorem~\ref{thm:classical} is then the textbook statement for distinct points. B\'ezout's theorem
itself is not formalised here.
\end{remark}

\section{Finite fields}\label{sec:finite}

Table~\ref{tab:finite} records where each statement starts to hold. Here $q=|K|$.

\begin{table}[h]
\centering\small
\begin{tabular}{@{}>{\raggedright\arraybackslash}p{0.40\textwidth}
  >{\raggedright\arraybackslash}p{0.30\textwidth}>{\raggedright\arraybackslash}p{0.24\textwidth}@{}}
\toprule
Statement & Holds for & Fails for \\ \midrule
Line lemma (cubics) & $q\ge3$ & $q=2$ (formal counterexample)\\
Conic line lemma, \decl{conic_three} & all $q$ & --- \\
Eight-point criterion, Theorems~\ref{thm:gp} and~\ref{thm:ncf} & all $q$ (vacuous for $q=2$) & --- \\
Theorem~\ref{thm:classical} (exactly nine) & $q=2$ (vacuous), $q=3$, $q\ge8$; all formalised & $q=4,5,7$\\
Pascal, Pappus via Cayley--Bacharach & $q\ge 11$ (\decl{CardGT}~\lean{K 10}) & not investigated \\
Pascal, Pappus, direct proofs (Section~\ref{sec:l1415}) & all $q$ & --- \\
\bottomrule
\end{tabular}
\caption{Field-size thresholds.}
\label{tab:finite}
\end{table}

Theorem~\ref{thm:classical} does not simply start to hold at some threshold. It holds
vacuously for $q=2$, holds for $q=3$, fails for $q=4,5,7$, and holds again from $q=8$ on. The
reason is that every counterexample needs a common component, since otherwise
Theorem~\ref{thm:ncf} applies. For small $q$ a line or conic has few points, and a pencil with a
common component can then have exactly nine rational common zeros. For $q\ge8$ the line or conic
alone has at least nine points, and the proof above goes through. For $q=2$ there are only seven
points. For $q=3$ lines and conics are too small even for that:

\begin{theorem}[$q=3$]\label{thm:f3}
Let $|K|=3$. If $F,G$ are linearly independent cubic forms with nine pairwise non-proportional
common zeros $P_1,\dots,P_9$ (not necessarily all their common zeros), then every cubic vanishing at
$P_1,\dots,P_8$ vanishes at $P_9$ (\decl{cayley_bacharach_card3}). Consequently the textbook form
holds over a finite field $K$ whenever $|K|\notin\{4,5,7\}$
(\decl{cayley_bacharach_classical_finite}).
\end{theorem}

\begin{proof}[Proof sketch]
We show that $F$ and $G$ have no common component, and then apply
\decl{cayley_bacharach_of_no_common_factor}. \emph{Common line $\ell$}
(\decl{no_common_line_card3}): at most $q+1=4$ of the nine points lie on $\ell$, so at least five
lie off $\ell$ and are common zeros of the two residual conics. No four of these are collinear: a
line through four of them also meets $\ell$, giving five points on a line in $\PP^2(\F_3)$, which has
only four. Conics through five points with no four collinear form a space of dimension at most
one, so the residual conics are proportional, and then so are $F$ and $G$. \emph{Common conic $q$
without a line component} (\decl{no_common_conic_card3}): at most one of the nine points is off
$q$, since two would force a common line. So $q$ has eight pairwise non-proportional zeros. But
projecting from one zero, the others lie on distinct lines through it (three collinear zeros would
give a line component), so $q$ has at most $q+2=5$ zeros. Both cases are contradictory.
The consequence combines this with \decl{cayley_bacharach_classical_card8}, the vacuous case
$q=2$, and the fact that no field has six elements.
\end{proof}

The proof is a counting argument; no case enumeration is done in Lean. Earlier versions of this
paper reported the case $q=3$ only as the result of an exhaustive computer search, which agrees:
pencils with a common line have at most eight common zeros over $\F_3$, and pencils with a common
conic having no line component at most five. The theorem is not vacuous over $\F_3$: $x^3-xz^2$ and
$y^3-yz^2$ have exactly the nine affine points as common zeros.

\begin{example}[counterexamples]\label{ex:counter}
In each case $F$ and $G$ are linearly independent, their common zeros in $\PP^2(\F_q)$ are
exactly the nine listed points, and $H$ vanishes at eight of them but not at the remaining one
(which plays the role of $P_9$). We write
$\F_4=\F_2[w]/(w^2+w+1)$. The common zeros and the values of~$H$ were computed by exhaustive
evaluation over $\PP^2(\F_q)$ with independent scripts.
\begin{itemize}
\item $q=4$: $F=z(x+wy+(w+1)z)(x+wy+wz)$ and $G=z(x+y+(w+1)z)(x+y)$. The common zeros are
$(0:1:0)$, $(1:0:0)$, $(1:0:w)$, $(1:1:0)$, $(1:1:1)$, $(1:1:w)$, $(1:w:0)$, $(1:w:1)$ and
$(1:w{+}1:0)$. The cubic $H=x^2z+w\,xyz+(w+1)\,y^2z$ vanishes at all of them except $(1:0:w)$.
\item $q=5$: $F=z(x+2y+3z)(x+2z)$ and $G=z(x+3y+2z)(x+y)$. The common zeros are $(0:1:0)$,
$(0:1:1)$, $(1:0:0)$, $(1:0:2)$, $(1:1:0)$, $(1:2:0)$, $(1:3:0)$, $(1:4:0)$ and $(1:4:2)$. The
cubic $H=x^2z+z^3$ vanishes at all of them except $(0:1:1)$.
\item $q=7$: $F=z(x+4y+4z)^2$ and $G=z(x+4y+z)(y+2z)$. The common zeros are the eight points of
$z=0$ and $(1:3:2)$. The cubic $H=x^2z$ vanishes at all of them except $(1:3:2)$.
\item $q=7$, common conic: $c=xz+y^2$, $F=z\,c$ and $G=x\,c$. The common zeros are the eight points
of $c$ and $(0:1:0)$. The cubic $H=y\,c$ vanishes at all of them except $(0:1:0)$.
\end{itemize}
These examples do not contradict the formal theorems. In each of them five or more of the first
eight points are collinear, or the first eight lie on a conic, and $F$ and $G$ share a component.
% Script: code/counterexamples.py in the repository.
\end{example}

The counterexamples are not formalised in Lean. The finite-field criterion and the counts behind
\decl{conic_points} were also checked by brute force for small $q$. The criterion was tested on
all $8$-subsets of $\PP^2(\F_3)$ and $\PP^2(\F_4)$ and on $20{,}000$ random $8$-subsets for $q=5$
and $q=7$, with no failures. Every conic over $\F_q$ with at least two zeros and no line
component was found to have exactly $q+1$ zeros, exhaustively for $q=3,4,5,7$ and on $40{,}000$
random conics for $q=8,9$.

\section{Applications}\label{sec:apps}

\subsection{Pascal and Pappus}

Let $A_1,\dots,A_6$ be six points and $L_{ij}=A_i\times A_j$ the side lines. The cubics
$F=L_{12}L_{34}L_{56}$ and $G=L_{23}L_{45}L_{61}$ both vanish at the six vertices and at
$X=L_{12}\cap L_{45}$, $Y=L_{23}\cap L_{56}$ and $Z=L_{34}\cap L_{61}$. If the vertices lie on a
conic~$q$, the cubic $H=(X\times Y)\cdot q$ vanishes at the six vertices and at $X$ and~$Y$. By
Cayley--Bacharach it vanishes at~$Z$. If $q(Z)\ne0$, then $Z$ lies on the line $XY$.

The basic formal statements \decl{pascal_of_cayley_bacharach} and
\decl{pappus_of_cayley_bacharach} carry the hypotheses of Theorem~\ref{thm:classical}
explicitly. These are that the nine points are distinct, that each side of $F$ differs from each
side of~$G$, and that $Z$ is off the conic ($q(Z)\ne0$, or $l\cdot Z\ne0\ne m\cdot Z$ for Pappus). The proofs discharge \lean{ExactlyMeetIn} and linear
independence of $F$ and $G$. We then derive all these hypotheses from the natural ones.

\begin{theorem}[Pascal]\label{thm:pascal}
Let $|K|\ge11$ or $K$ infinite. Let $A_1,\dots,A_6$ be distinct points on a conic $q$ that is not
a product of two linear forms over~$K$ (in particular $q\ne0$). Then the points $L_{12}\times L_{45}$,
$L_{23}\times L_{56}$ and $L_{34}\times L_{61}$ are collinear.
\end{theorem}

\begin{lstlisting}
theorem pascal_clean [CardGT K 10] (A₁ A₂ A₃ A₄ A₅ A₆ : Fin 3 → K)
    (hA : ∀ i j, i ≠ j → ![A₁, A₂, A₃, A₄, A₅, A₆] i ⨯₃ ![A₁, A₂, A₃, A₄, A₅, A₆] j ≠ 0)
    (q : Fin 6 → K) (hirr : ∀ ℓ m : Fin 3 → K, q ≠ linMul ℓ m)
    (hq₁ : evalC q A₁ = 0) (hq₂ : evalC q A₂ = 0) (hq₃ : evalC q A₃ = 0)
    (hq₄ : evalC q A₄ = 0) (hq₅ : evalC q A₅ = 0) (hq₆ : evalC q A₆ = 0) :
    Matrix.det ![A₁ ⨯₃ A₂ ⨯₃ (A₄ ⨯₃ A₅), A₂ ⨯₃ A₃ ⨯₃ (A₅ ⨯₃ A₆),
      A₃ ⨯₃ A₄ ⨯₃ (A₆ ⨯₃ A₁)] = 0
\end{lstlisting}

The conclusion $\det(X,Y,Z)=0$ is the usual coordinate test for collinearity. The genericity
conditions come from \decl{pascal_generic}. First, no three vertices are collinear
(\decl{noncollinear_of_irreducible}, via \decl{conic_three}). From this we get the nine
side-line conditions, $X,Y,Z\ne0$, that $X,Y,Z$ are distinct from each vertex and from each
other, and that $X,Y,Z$ are off the conic (\decl{pascal_points_off_conic}). Irreducibility is used to establish the genericity this proof needs; it is not needed for the
conclusion itself, and the determinant identity $\det(X,Y,Z)=0$ holds for degenerate conics too~\cite{PascalHexagon}. For the line pair $q=xy$, with $A_1,A_2,A_3$ on $x=0$ and
$A_4,A_5,A_6$ on $y=0$, we get $L_{12}=L_{23}$ and $X=Y$. With $A_1,A_2,A_4,A_5$ on $x=0$ we get
$L_{12}=L_{45}$ and $X=0$. In both cases $\det(X,Y,Z)=0$ still holds, but the nine points are no
longer distinct, so Cayley--Bacharach does not apply. Pappus's theorem is the alternating-vertex
specialisation to a line pair.

\begin{theorem}[Pappus]\label{thm:pappus}
Let $|K|\ge11$ or $K$ infinite. Let $A_1,A_3,A_5$ be distinct points on a line~$l$ and
$A_2,A_4,A_6$ distinct points on a line~$m$, none of the six lying on both lines. Then the points
$L_{12}\times L_{45}$, $L_{23}\times L_{56}$ and $L_{34}\times L_{61}$ are collinear.
\end{theorem}

This is \decl{pappus_clean}. The condition ``none on both lines'' is stated as
$m\cdot A_1\ne 0$, $l\cdot A_2\ne0$, and so on. The hypotheses imply $l\ne m$
(\decl{pappus_lines_distinct}), so this need not be assumed. \decl{pappus_generic} derives the
genericity hypotheses. In the tested collisions and line-intersection degenerations (\texttt{check\_stage8.py}), the
nine-point genericity fails, although $\det(X,Y,Z)=0$ still holds: two Pascal points coincide or
vanish, or a Pascal point becomes a vertex. These are therefore
limits of the proof method, not counterexamples to Pappus's theorem. A direct polynomial proof
would not need them. A formalisation of Pascal's and Pappus's theorems over arbitrary fields by
such a direct proof exists~\cite{PascalHexagon}, and version~5 of this development contains direct
proofs over every field as well (Section~\ref{sec:l1415}). The point of the present section is that both theorems follow
from Cayley--Bacharach.

\subsection{Line pairs and the Gram determinant}

Pascal's theorem is often stated for a nondegenerate conic, $\det M(q)\ne0$, where $M(q)$ is the
symmetric matrix with $v^{\mathsf T}M(q)v=q(v)$. The formal \decl{gram}~$q$ has the off-diagonal
coefficients halved, and \decl{gram_quadForm} shows $v^{\mathsf T}M(q)v=q(v)$ when $2\ne0$ in~$K$.

\begin{proposition}\label{prop:gram}
Let $2\ne0$ in $K$. If $q=\ell\cdot m$, then $\det M(q)=0$. Hence, for $|K|>10$, Pascal's theorem
holds for conics with $\det M(q)\ne 0$. If $K$ is algebraically closed, then conversely $\det M(q)=0$
implies that $q$ is a product of two linear forms.
\end{proposition}

The three statements are \decl{det_gram_linMul}, \decl{pascal_clean_det} and
\decl{det_gram_eq_zero_iff}.

For the converse, a kernel vector $v$ of $M(q)$ is a zero of $q$ at which the polar form
$B(v,\cdot)$ vanishes. Over an algebraically closed field, $q$ has a further zero $w$ not
proportional to~$v$. Then $q(sv+tw)=s^2q(v)+st\,B(v,w)+t^2q(w)=0$, with $B$ the unhalved polar form as above, so $q$ contains the line $vw$,
and the conic line lemma splits it off. The hypothesis $2\ne0$ is needed in the forward
direction. In Lean $2^{-1}=0$ in characteristic~2, so \decl{gram} is not the Gram matrix there.
\decl{det_gram_linMul_char_two} records that the square $(x+y+z)^2$ over $\mathbb{Z}/2$ then gets
determinant~1. As predicates on conics, $\det M(q)\ne0$ is strictly stronger than irreducibility over a field
that is not algebraically closed: $x^2+y^2$ over $\mathbb{Q}$ is not a product of linear forms but
has $\det M=0$. This does not give more six-point configurations, however: that conic has only the
rational point $(0:0:1)$, and in characteristic $\ne2$ a singular conic with two distinct rational
zeros splits. The genuine additional scope of \decl{pascal_clean} is characteristic two.

\section{Does the statement mean what it says?}\label{sec:correct}

The development uses its own coefficient-vector representation. A reader has to check that the
basic operations mean what their names suggest. These lemmas are provided for that purpose.
\begin{itemize}
\item \decl{eval_expand} writes out \lean{eval F v} as the explicit cubic polynomial in
$v_0,v_1,v_2$. \decl{evalC_expand} does the same for conics.
\item \decl{eval_mulLin} and \decl{evalC_linMul} show that \lean{mulLin} and \lean{linMul} are
the products they are meant to be: $(\ell q)(v)=(\ell\cdot v)\,q(v)$ and
$(\ell m)(v)=(\ell\cdot v)(m\cdot v)$.
\item \decl{eval_smul_vec} gives $F(cv)=c^3F(v)$ and \decl{vanishesAt_smul_iff} shows that
vanishing at a point does not depend on the representative. \decl{evalC_smul_vec} gives
$q(cv)=c^2q(v)$.
\item \decl{ker_evalMap_eq} identifies the kernel of the evaluation map with the cubics
vanishing at the points. The dimensions in Theorem~\ref{thm:criterion} are those of this kernel.
\item The two easy directions of Theorem~\ref{thm:criterion} show that neither \lean{h5} nor
\lean{hconic} can be dropped from the criterion: without either one the dimension can be~3. The
general-position form of Cayley--Bacharach then fails: for five points on $y=0$ plus three more, our numerical check exhibits two cubics through
the eight points and a ninth point~$R$, and a third cubic through the eight that misses~$R$.
\item The hypothesis \lean{hP} (pairwise non-proportional) is exactly distinctness in
$\PP^2(K)$. \decl{card_le_seven_of_not_atLeastThree} shows that it cannot hold for eight points
over~$\F_2$, so the \lean{\_any} theorems are vacuous there. This is the expected behaviour,
since $\PP^2(\F_2)$ has seven points.
\item \lean{ExactlyMeetIn} is literally ``the common zeros are the $P_i$ up to scaling''. The
hypothesis refers to rational points; see the remark after Theorem~\ref{thm:classical}.
\item The three sharpness statements \decl{line_lemma_false_zmod_two}, \decl{det_gram_linMul_char_two}
and the $\F_4,\F_5,\F_7$ counterexamples show that the size and characteristic hypotheses are
not artefacts of the proofs. The first two are formal; the last is computational.
\end{itemize}
A further file, \texttt{Bridge.lean}, records that the local definitions are the standard
notions. \decl{eval_smul_smul} and \decl{evalC_smul_smul} give $(aF)(cv)=ac^3F(v)$ and
$(aq)(cv)=ac^2q(v)$, and \decl{vanishesAt_smul_smul_iff} shows that \decl{VanishesAt} is well defined
on forms and points up to nonzero scalars; \decl{vanishesAt_iff}, \decl{commonZero_iff} and
\decl{exactlyMeetIn_iff} (with \decl{cross_eq_zero_iff_eq_smul}) read \decl{VanishesAt},
\decl{CommonZero} and \decl{ExactlyMeetIn} through \lean{eval}. \decl{mem_onLine_iff_mem_ker} says
that \decl{onLine} is the kernel of the linear form~$\ell$. \decl{col4_iff_exists_mulVec},
\decl{col4_iff_rank_le_two} and \decl{col4_iff_minors} show that \decl{Col4} holds iff the
$4\times3$ coordinate matrix has rank at most~2, iff all its $3\times3$ minors vanish.
\decl{polarC_eq}, \decl{polarC_self}, \decl{polarC_comm} and \decl{polarC_eq_bilin} show that
\decl{polarC} is the symmetric bilinear form $q(x+y)-q(x)-q(y)$, with no division, in every
characteristic. These lemmas are not used elsewhere and add no new results.

The conclusions of Pascal and Pappus are stated as $\det(X,Y,Z)=0$. We do not connect this here
with Mathlib's \lean{Projectivization.IsCollinear}. That connection is proved in the author's
Pascal development~\cite{PascalHexagon}.

\section{Related work}\label{sec:related}

\paragraph{Mathematics.} The classical sources are Cayley's paper of 1843~\cite{Cayley} and
Bacharach's paper of 1886~\cite{Bacharach}. Eisenbud, Green and Harris~\cite{EGH} trace the
history of the theorem and its later generalisations. Fulton's textbook
proves Pascal's and Pappus's theorems from Max Noether's fundamental theorem, over an
algebraically closed field~\cite[\S5.6]{Fulton}. Our proof is the elementary linear-algebra
argument. We have not found the finite-field thresholds of Section~\ref{sec:finite} stated in the
literature, but we did not search the literature on curves over finite fields systematically,
and we make no claim that they are new.

\paragraph{Formalisations.} We searched for existing formalisations in the following places. For
GitHub we used code and repository search and the list of open Mathlib pull requests. We also
checked the Isabelle Archive of Formal Proofs (1{,}034 entries at the time)~\cite{AFP}, the HOL~Light
repository~\cite{HOLLight}, and the list of famous theorems proved in Rocq~\cite{Rocq100}. The
searches were made between 28 and 30 September 2026 and found no formalisation of the theorem.
% TODO(Josh): Mizar MML not searched; search it or keep this sentence.
We did not search the Mizar Mathematical Library.
Mathlib's list of 1000 theorems has an entry ``Cayley--Bacharach theorem'' with no associated
declaration~\cite{Mathlib1000}. The following nearby items exist.
\begin{itemize}
\item Google DeepMind's \emph{formal-conjectures} states the Eisenbud--Green--Harris conjecture,
a much more general open problem motivated by Cayley--Bacharach. It does not prove
Cayley--Bacharach~\cite{FormalConjectures}.
\item In Meta's \emph{atlas-lean}, one file states B\'ezout's theorem and Pascal's theorem. The
B\'ezout statements and a step of the Pascal construction are left as \lean{sorry}. A GitHub code search of
that repository for ``Bacharach'' returned no results~\cite{AtlasLean}.
\item The \emph{ShadowBench} autoformalisation benchmark includes an informal statement and proof of
the theorem as a statement-formalisation task. No formal proof is public in that item; the benchmark's gold
formalisations are private and we did not inspect them~\cite{ShadowBench}.
\end{itemize}
A search of the Lean Zulip chat for ``Bacharach'' on 30 September 2026 returned no results. As far as we could determine, this is the first
formalisation of the Cayley--Bacharach theorem in a proof assistant, but the searches are not
exhaustive.

\section{Verification}\label{sec:gate}

The library has 19 modules and a root file, about 4{,}150 lines, 216 theorems and 32
definitions (at the published commit; version~5 adds the two modules of Section~\ref{sec:l1415}). At the published commit, a full \lean{lake build} finishes with 0 errors and 0 warnings. Lean reports every
use of \lean{sorry} as a warning, so this rules it out. A grep of all sources finds no
\lean{sorry}, \lean{admit}, \lean{native\_decide}, \lean{axiom}, \lean{set\_option},
\lean{\#eval}, \lean{run\_cmd}, \lean{macro}, \lean{elab}, \lean{syntax}, \lean{notation} or
\lean{import Lean}. On the published-commit sources, a single scratch file outside the library
(\path{scratch/PaperAxioms.lean}, output \path{scratch/paper_axioms.out}) runs \lean{\#print axioms}
on the 99 declarations cited in the paper before version~5 and one further supporting theorem. Every reported
axiom is among \lean{propext}, \lean{Classical.choice} and \lean{Quot.sound}; the abbreviation
\decl{PlaneCubic} uses none. The repository's \path{gate.sh} regenerates these queries itself for the same 100 fully qualified
declarations, rejects forbidden tokens and fails on any other axiom; on a fresh clone of commit
\texttt{\shortcommit} it prints \texttt{PASS (100 declarations, standard axioms only)} with GNU
tools (on macOS the published script stops at a GNU-only \lean{mktemp} option; see
Section~\ref{sec:l1415}).
\path{tests/tamper.sh} plants six kinds of defect in scratch copies (a \lean{sorry}, a split
\lean{axiom}, a fake \lean{\#eval} report, a macro, a weakened field-size hypothesis, a wrapped
extra axiom) and requires each to be rejected with a nonzero exit status and the expected failure
message; all six are rejected.

Python scripts in \texttt{code/} check the mathematics independently of Lean.
\begin{itemize}
\item \texttt{check\_cb.py} (exact rationals, sympy). For five random configurations of eight
points it verifies $\dim V(P)=2$, computes the ninth base point, and checks that an independently
computed cubic through the eight vanishes there. It also checks the five-collinear negative
example of Section~\ref{sec:correct}.
\item \texttt{check\_stage5.py}. This checks the $\mathbb{R}$ example $x q, y q$ with
$q=x^2+y^2+z^2$, which has the single real common zero $(0:0:1)$, and the projection construction
of \decl{false_of_conic}.
\item \texttt{check\_stage6.py}, \texttt{check\_stage7.py}, \texttt{check\_stage8.py},
\texttt{examples\_stage8.py}. These test random Pascal and Pappus configurations and the
degenerations of Section~\ref{sec:apps}. In a negative control, moving a vertex off the conic gave a
nonzero determinant in 50 of 50 observed cases.
\item \texttt{criterion\_examples.py}, \texttt{counterexamples.py}, and five search scripts
(\texttt{brute9.py}, \texttt{classical9.py}, \texttt{classical9b.py}, \texttt{q8conic.py},
\texttt{count10.py}). These contain the rank computations of Section~\ref{sec:correct}, the brute-force checks over $\F_q$ and
the counterexamples of Example~\ref{ex:counter}; saved outputs of the long searches are in
\texttt{code/outputs/}.
\end{itemize}
Automated checks cannot replace comparing the formal statements with their informal
counterparts. Sections~\ref{sec:conv} and~\ref{sec:correct} are meant to make that comparison
easier.
The author has checked each statement in Table~\ref{tab:thms} against its informal counterpart by hand.
The Layer-14 and Layer-15 statements of Section~\ref{sec:l1415} have not been checked by hand.

\subsection{Layers 14 and 15}\label{sec:l1415}

Version~5 adds two further checks, in the files \path{CayleyBacharach/L14.lean} (1{,}050 lines,
namespace \lean{CBL14}) and \path{CayleyBacharach/L15Classical.lean} (1{,}040 lines, namespace
\lean{CayleyBacharach.Classical}), at the new commit \texttt{d9c674c68c0a1b475a284f9bd33ec5f7e6d52689}. The gate explicitly lists 27 Layer-14
declarations and 41 Layer-15 declarations in \path{gate.sh}, which prints \texttt{PASS (168 declarations, standard axioms only)}
(100 earlier, 27 from Layer~14, 41 from Layer~15), and \path{tests/tamper.sh} still rejects all six
planted defects. The earlier library files are byte-for-byte unchanged. The published \path{gate.sh}
used GNU \lean{mktemp} options and failed on macOS (BSD tools) before checking anything. The new commit
uses a portable temporary file name, and the gate passes with BSD tools. \path{tests/tamper.sh} still
needs GNU \lean{sed}.

\paragraph{Layer 14: statements in Mathlib's language, with witnesses.}
In \path{CayleyBacharach/L14.lean} the main results are restated with Mathlib-only types: forms are
homogeneous \lean{MvPolynomial (Fin 3) K}, points are vectors \lean{Fin 3} $\to$ \lean{K}, distinctness is
\lean{LinearIndependent K ![u, v]}, vanishing is \lean{MvPolynomial.eval}, common components are
factorisations in \lean{MvPolynomial}, the cubics through the points form
\lean{homogeneousSubmodule} intersected with evaluation kernels, and collinearity is
$\det(X,Y,Z)=0$. Each restatement is proved from the library theorem through bridge lemmas
(\decl{l14b_eval3}, \decl{l14b_cubicsThrough}, \decl{l14b_lin}, \decl{l14b_con}, \decl{l14b_exact} and
others). Twelve restatements cover nine results: Theorems~\ref{thm:gp} and~\ref{thm:ncf}
(\decl{l14_cayley_bacharach_any}, \decl{l14_cayley_bacharach_of_no_common_factor_any});
Theorem~\ref{thm:classical} (\decl{l14_cayley_bacharach_classical}, and
\decl{l14_cayley_bacharach_classical_finite} for $|K|\notin\{4,5,7\}$, which subsumes the $|K|\ge8$ and
$|K|=3$ cases); Theorem~\ref{thm:f3} (\decl{l14_cayley_bacharach_card3}); Theorem~\ref{thm:criterion}
(\decl{l14_finrank_eq_two_iff}); Theorems~\ref{thm:pascal} and~\ref{thm:pappus} (\decl{l14_pascal},
\decl{l14_pappus}, and \decl{l14_pascal_points}, \decl{l14_pappus_points}, in which $X,Y,Z$ are
\emph{any} points satisfying their two incidences); the cubic part of Lemma~\ref{lem:line}
(\decl{l14_line_lemma}); and \decl{line_lemma_false_zmod_two} (\decl{l14_line_lemma_false_zmod_two}).
A mechanical check of the elaborated types found no constant of this library in any of them. No
restatement turned out false or needed an extra hypothesis. Six witness theorems prove every hypothesis
on an explicit configuration, together with non-degeneracy: the $3\times3$ grid over $\mathbb{Q}$, with
the nine points pairwise distinct, $F,G$ independent, $H=F+G$ proportional to neither, exactness and the
absence of common components proved (\decl{l14w_cayley_bacharach}); the same pencil over $\F_3$, whose
common zeros are exactly the nine affine points (\decl{l14w_cayley_bacharach_zmod3}); both directions of
the criterion, eight grid points with $\dim=2$ and eight points with five on $y=0$ with $\dim\ne2$
(\decl{l14w_finrank_eq_two_iff}); six rational points on $x^2+y^2=z^2$, with the conic proved not a line
pair and the three Pascal points pairwise distinct (\decl{l14w_pascal}); a Pappus configuration over
$\mathbb{Q}$ with pairwise distinct Pappus points (\decl{l14w_pappus}); and $F=z(x^2+y^2+z^2)$ for the
line lemma (\decl{l14w_line_lemma}). The $\F_2$ counterexample is itself concrete. Not restated, to stay
within the time box: Proposition~\ref{prop:gram} (\decl{det_gram_eq_zero_iff}, \decl{det_gram_linMul},
\decl{pascal_clean_det}), whose Gram matrix is a local definition, the conic part of
Lemma~\ref{lem:line}, and the auxiliary lemmas (\decl{conic_three}, \decl{cubic_four} and others).

\paragraph{Layer 15: comparison with the classical statements.}
Another agent collected textbook statements of the four topics from the literature, with citations
(Eisenbud--Green--Harris~\cite{EGH}, Fulton~\cite[\S5.6]{Fulton}, Reid's \emph{Undergraduate
Algebraic Geometry}, Coxeter's \emph{Projective Geometry}, Wikipedia); page numbers were given as best
known. An AI model given only these statements, with no access to our code or to this paper, wrote them
in Lean with Mathlib notions (points of Mathlib's \lean{Projectivization}, forms in \lean{MvPolynomial},
\lean{IsRelPrime} for ``no common component'') and definitions of its own, among them an intersection
multiplicity defined as a module length over a localisation. It wrote 11 theorems and 4 well-formedness
obligations, all with \lean{sorry}. Its transcript and self-reported file list were audited. A comparer
preserved every binder and conclusion, mechanically rewriting each theorem with a placeholder proof
as a Prop-valued definition. It proved coverage by library derivations or new proofs, or recorded a
narrower or out-of-scope result, without changing the existing library files. Of the 15 statements, 9 are \emph{covered} (proved), 3 are covered only in a
\emph{narrower} sense, 3 are \emph{out of scope} for this paper (2 of them nevertheless proved), and none
was classified as \emph{false as written}. The proofs are collected, with the writer's definitions and mechanically
rewritten statement definitions preserved from the comparison file, in \path{CayleyBacharach/L15Classical.lean}, namespace
\lean{CayleyBacharach.Classical}. All of them depend only on the standard axioms. The covered statements
are Cayley--Bacharach for coprime cubics with nine distinct common points, over every field, in both
the ``passes through $P_9$'' and the ``$H=aF+bG$'' forms (\decl{cov_CB1_a}, \decl{cov_CB1_a_linear}); the
criterion under ``no four collinear, no seven on a conic'' (\decl{cov_CB2_a}); Pascal, Pappus and
Pappus in hexagon form; and the three corresponding obligations, that opposite sides meet in a unique
point.

\emph{New proofs} (not in the main files):
\begin{itemize}
\item \decl{cov_CB3_a}: Pascal's theorem over every field, for an irreducible conic and six distinct
points, with $X,Y,Z$ given by their incidences (the library's \decl{pascal_clean} needs $|K|>10$). The
explicit conic $Q(p)=[X,\,BC\cap Ep,\,CD\cap pA]$ passes through $A,\dots,E$ and is nonzero, so it is
proportional to the given conic by \decl{conics_le_one}, and $Q(F)=0$ is the collinearity of $X,Y,Z$. This
proof does not go through Cayley--Bacharach.
\item \decl{cov_CB3_b}: the Braikenridge--Maclaurin converse over every field: if no three of six points
are collinear and $X,Y,Z$ are collinear, the six points lie on a conic. The main files have no converse.
\item \decl{cov_CB4_a}, \decl{cov_CB4_b}: Pappus's theorem over every field, in the three-point form and
in hexagon form (the library's \decl{pappus_clean} needs $|K|>10$). The key step \decl{l15_papid} is a
single polynomial identity.
\item \decl{l15_intMult_ne_zero_iff}: the writer's intersection multiplicity of two forms at a point is
nonzero if and only if both forms vanish there. The comparer did not verify that this definition
agrees with Fulton's local intersection multiplicity; only its support equivalence was proved.
\end{itemize}
The comparer also proved Pascal and both Pappus forms through the library route, with the extra
hypothesis $|K|>10$ (\decl{cov_CB3_a_part}, \decl{cov_CB4_a_part}, \decl{cov_CB4_b_part}).

\emph{Narrower} (proved part, and what is missing):
\begin{itemize}
\item The writer's proposed encoding of Fulton's multiplicity version, with simple points or with $C_1$ irreducible (two statements):
proved when $P_1,\dots,P_9$ are distinct, without algebraic closure (\decl{cov_CB1_b_part},
\decl{cov_CB1_b_irreducible_part}); missing is the case of coincident points, that is genuine
intersection multiplicities, which needs B\'ezout's theorem and local intersection numbers.
\item The ninth base point of the pencil when two members meet in nine distinct points: existence over
every field (\decl{cov_CB2_b_distinct_part}); missing is uniqueness, that two coprime cubics have at most
nine common rational points (B\'ezout).
\end{itemize}

\emph{Out of scope}: the ninth base point over an algebraically closed field as the residual point of the
intersection cycle (not proved; it needs B\'ezout and intersection cycles), and Braikenridge--Maclaurin
with its obligation (proved, as above).

\emph{Writer statements found false as written}: none; this does not establish the truth of the
unproved statements.

\section{Limitations and further work}

\begin{itemize}
\item \emph{Multiplicities.} The nine points are distinct. The version with intersection
multiplicities, which covers tangencies and degenerate hexagons, needs B\'ezout's theorem and
local intersection numbers. Neither is formalised here. The Layer-15 file proves only the case of
distinct points of the writer's proposed encoding of Fulton's multiplicity version (Section~\ref{sec:l1415}).
\item \emph{Rational points.} Theorem~\ref{thm:classical} is about $K$-rational common zeros.
We do not prove that ``no common component'' implies ``exactly nine points over $\bar K$, counted
with multiplicity'', nor that two coprime cubics have at most nine common rational points; the
uniqueness half of the corresponding Layer-15 statement is therefore open.
\item \emph{Counterexamples over $\F_4$, $\F_5$, $\F_7$.} These are verified by computer, not in
Lean, so the ``only if'' half of the finite-field classification is not formalised.
\item \emph{Pascal and Pappus over small fields.} The versions derived from Cayley--Bacharach carry
\decl{CardGT}~\lean{K 10} because they go through the $|K|\ge11$ version of
Theorem~\ref{thm:classical}. Routing them through \decl{cayley_bacharach_classical_card8} should lower
this to $|K|\ge8$, and proving directly that the two triangle cubics have no common component might
remove the size hypothesis from that route. Neither has been done. Pascal and Pappus themselves now
hold over every field by the direct proofs of version~5 (Section~\ref{sec:l1415}).
\item \emph{Representation.} Cubics and conics are coefficient vectors, not Mathlib's
\lean{MvPolynomial}, and points are vectors, not elements of \lean{Projectivization}. This keeps
the linear algebra simple. A contribution to Mathlib would need a bridge to the standard
structures. The Layer-14 file bridges forms to \lean{MvPolynomial} for the restated theorems, and the
Layer-15 file connects some of them to points of \lean{Projectivization}
(Section~\ref{sec:l1415}); the main files are unchanged.
\item \emph{Further applications.} Associativity of the group law on a plane cubic is the best-known
application of Cayley--Bacharach. It needs the chord--tangent construction and multiplicities, and
it is not attempted.
\end{itemize}

\section*{Acknowledgements and use of AI}
A large-language-model AI assistant was used to write the Lean proofs and the Python checking
scripts, to run the prior-art searches, and to draft this paper. The author directed the work and
reviewed the paper, and checked each formal statement in Table~\ref{tab:thms} against its
informal counterpart by hand.
The Layer-14 and Layer-15 files, the collection of classical statements, the blind formalisation and
the comparison (Section~\ref{sec:l1415}) were all done by AI agents and have not been checked by hand.
Every proof is checked by the Lean kernel, and the author takes responsibility for all content.

\begin{thebibliography}{99}\sloppy
\bibitem{AFP} \emph{Archive of Formal Proofs}, \url{https://www.isa-afp.org/}, accessed 28
September 2026.
\bibitem{AtlasLean} Meta Platforms, Inc., \emph{atlas-lean}, file
\path{v1/Atlas/AlgebraicGeometryI/code/Lec5BezoutPascal.lean},
\url{https://github.com/facebookresearch/atlas-lean}, accessed 30 September 2026.
\bibitem{Bacharach} I.~Bacharach, Ueber den Cayley'schen Schnittpunktsatz, \emph{Mathematische
Annalen} 26 (1886), 275--299.
\bibitem{PascalHexagon} J.~Bald, \emph{pascal-hexagon}: Pascal's hexagon theorem over an arbitrary
field in Lean~4, \url{https://github.com/jpbald93/pascal-hexagon}, accessed 30 September 2026.
\bibitem{Cayley} A.~Cayley, On the intersection of curves, \emph{Cambridge Mathematical Journal}
3 (1843), 211--213.
\bibitem{Lean4} L.~de Moura and S.~Ullrich, The Lean~4 theorem prover and programming language, in
\emph{Automated Deduction -- CADE~28}, Lecture Notes in Computer Science 12699, Springer, 2021,
625--635.
\bibitem{EGH} D.~Eisenbud, M.~Green and J.~Harris, Cayley--Bacharach theorems and conjectures,
\emph{Bulletin of the American Mathematical Society (N.S.)} 33 (1996), no.~3, 295--324.
\bibitem{FormalConjectures} Google DeepMind, \emph{formal-conjectures}, file
\path{FormalConjectures/Paper/EisenbudGreenHarris.lean},
\url{https://github.com/google-deepmind/formal-conjectures}, accessed 30 September 2026.
\bibitem{Fulton} W.~Fulton, \emph{Algebraic Curves: An Introduction to Algebraic Geometry},
first published 1969; edition of January 2008, made freely available by the author.
\bibitem{HOLLight} J.~Harrison, HOL~Light repository, \url{https://github.com/jrh13/hol-light},
accessed 28 September 2026.
\bibitem{Mathlib} The mathlib Community, The Lean mathematical library, in \emph{Proceedings of the
9th ACM SIGPLAN International Conference on Certified Programs and Proofs (CPP 2020)}, ACM, 2020,
367--381.
\bibitem{Mathlib1000} The mathlib Community, \texttt{docs/1000.yaml}, Mathlib repository,
\url{https://github.com/leanprover-community/mathlib4}, accessed 30 September 2026.
\bibitem{Rocq100} \emph{Famous theorems proven in Rocq}, \url{https://madiot.fr/rocq100/}, accessed
30 September 2026.
\bibitem{ShadowBench} \emph{ShadowBench}, item
\path{data/algebraic-geometry/L3/alg_gen_L3_001}, \url{https://github.com/ldilab/shadowbench},
accessed 30 September 2026.
\bibitem{Wikipedia} Wikipedia, \emph{Cayley--Bacharach theorem},
\url{https://en.wikipedia.org/wiki/Cayley\%E2\%80\%93Bacharach_theorem}, accessed 30 September
2026.
\end{thebibliography}

\end{document}
